% !TeX root = ../main.tex
\section*{Обозначения и используемые факты}
\addcontentsline{toc}{section}{Обозначения и используемые факты}
Нейтральный элемент группы обозначаем $e$; в матричных группах это $I$, а в аддитивных группах --- $0$. Через $\Z_n$ обозначаем аддитивную
группу классов вычетов, через $\Z_n^*$ --- мультипликативную группу
обратимых классов. Пишем
\[
 U=\{z\in\C:|z|=1\},\qquad U_n=\{z\in\C:z^n=1\},
 \qquad \Rp=(0,+\infty).
\]
Операция в $U$, $U_n$, $\C^*$ и $\Rp$ --- умножение.
Диэдральная группа $D_n$ имеет порядок $2n$.
Произведение перестановок понимаем как композицию справа налево:
$(\sigma\tau)(i)=\sigma(\tau(i))$.
Левый смежный класс имеет вид $gH$, правый --- $Hg$;
индекс подгруппы записываем $[G:H]$.

\Def{Центр группы --- подгруппа
 $Z(G)=\{z\in G:zg=gz\text{ для всех }g\in G\}$.
 Нормализатор подгруппы $H$ ---
 $N_G(H)=\{g\in G:gHg^{-1}=H\}$.
 Централизатор элемента $x$ ---
 $C_G(x)=\{g\in G:gx=xg\}$.}

\lem{Порядок и степень}{
 Если $\ord(g)=m<\infty$, то $g^k=e$ тогда и только тогда,
 когда $m\mid k$. В частности,
 \[\ord(g^k)=\frac{m}{\gcd(m,k)}.\]
}{
 Разделим $k$ с остатком: $k=qm+r$, $0\le r<m$.
 Тогда $g^k=g^r$, а $g^r=e$ при $0<r<m$ противоречит
 минимальности $m$. Для второй формулы нужно найти наименьшее
 положительное $t$ с $m\mid kt$. Если $d=\gcd(m,k)$,
 условие равносильно $m/d\mid t$, поскольку $m/d$ и $k/d$
 взаимно просты.
}

\thm{Лагранжа}{
 Для конечной группы $G$ и подгруппы $H$ верно
 $|G|=[G:H]|H|$. Поэтому $|H|\mid |G|$ и $\ord(g)\mid |G|$.
}{
 Смежные классы разбивают $G$, а отображение $h\mapsto gh$
 является биекцией $H\to gH$. Для порядка элемента применяем
 формулу к циклической подгруппе $\langle g\rangle$.
}

\thm{О гомоморфизме}{
 Если $\varphi:G\to K$ --- гомоморфизм, то
 $\ker\varphi\trianglelefteq G$ и
 $G/\ker\varphi\cong\im\varphi$.
}{
 При $h\in\ker\varphi$ имеем
 $\varphi(ghg^{-1})=\varphi(g)e\varphi(g)^{-1}=e$.
 Отображение $g\ker\varphi\mapsto\varphi(g)$ корректно:
 два значения совпадают ровно тогда, когда $g_2^{-1}g_1\in\ker\varphi$.
 Оно сохраняет произведение и является биекцией на образ.
}

\thm{Об орбите и стабилизаторе}{
 При действии группы $G$ на множестве $X$ орбита $Gx$
 находится в биекции с множеством левых классов $G/G_x$,
 где $G_x=\{g:gx=x\}$. Если $G$ конечна, то
 $|Gx|=|G|/|G_x|$.
}{
 Биекция задаётся формулой $gG_x\mapsto gx$.
 Равенство $g_1x=g_2x$ равносильно $g_2^{-1}g_1\in G_x$,
 что одновременно доказывает корректность и инъективность.
}
